SEGLEA computes the source effective-address point and applies the selected segment\textquotesingle{}s pre-translation to that point. The size suffix selects the index scale used during effective-address calculation. On success, SEGLEA writes the segment result to the destination and clears FLAGS. A failed segment-bound check sets (:ref:base.registers.STATE.FLAGS.V:) while clearing (:ref:base.registers.STATE.FLAGS.Z:), (:ref:base.registers.STATE.FLAGS.N:), and (:ref:base.registers.STATE.FLAGS.C:), suppresses the destination write, commits effective-address auto-update effects, and completes without an exception. The B-sized 0 or 1 failure result supplies the REPcc observed value.

Let \(a\) be the source (:term:base.terms.effective_address:) before a memory read. The
effective-address form selects a (:term:base.terms.segment_register_image:), which supplies the
\(\mathit{base}\), \(\mathit{span}\), and \(\mathit{limit}\) quantities used by
(:term:base.terms.segment_pre_translation:).
A zero mantissa disables bounds and returns \(a\). In translated mode, \(a\) must satisfy
\(0 \le a < \mathit{span}\) and the result is \(\mathit{base}+a\). In bounds-only mode, \(a\) must satisfy
\(\mathit{base} \le a < \mathit{limit}\) and the result remains \(a\). An out-of-bounds address writes FLAGS as
(:ref:base.registers.STATE.FLAGS.Z:)\texttt{=}(:ref:base.registers.STATE.FLAGS.N:)\texttt{=}(:ref:base.registers.STATE.FLAGS.C:)\texttt{=0} and (:ref:base.registers.STATE.FLAGS.V:)\texttt{=1}, in which case the destination is not written. A successful address writes FLAGS as (:ref:base.registers.STATE.FLAGS.Z:)\texttt{=}(:ref:base.registers.STATE.FLAGS.N:)\texttt{=}(:ref:base.registers.STATE.FLAGS.C:)\texttt{=}(:ref:base.registers.STATE.FLAGS.V:)\texttt{=0}.
